% Options for packages loaded elsewhere
\PassOptionsToPackage{unicode}{hyperref}
\PassOptionsToPackage{hyphens}{url}
\documentclass[
]{article}
\usepackage{xcolor}
\usepackage[margin=1in]{geometry}
\usepackage{amsmath,amssymb}
\setcounter{secnumdepth}{-\maxdimen} % remove section numbering
\usepackage{iftex}
\ifPDFTeX
  \usepackage[T1]{fontenc}
  \usepackage[utf8]{inputenc}
  \usepackage{textcomp} % provide euro and other symbols
\else % if luatex or xetex
  \usepackage{unicode-math} % this also loads fontspec
  \defaultfontfeatures{Scale=MatchLowercase}
  \defaultfontfeatures[\rmfamily]{Ligatures=TeX,Scale=1}
\fi
\usepackage{lmodern}
\ifPDFTeX\else
  % xetex/luatex font selection
\fi
% Use upquote if available, for straight quotes in verbatim environments
\IfFileExists{upquote.sty}{\usepackage{upquote}}{}
\IfFileExists{microtype.sty}{% use microtype if available
  \usepackage[]{microtype}
  \UseMicrotypeSet[protrusion]{basicmath} % disable protrusion for tt fonts
}{}
\makeatletter
\@ifundefined{KOMAClassName}{% if non-KOMA class
  \IfFileExists{parskip.sty}{%
    \usepackage{parskip}
  }{% else
    \setlength{\parindent}{0pt}
    \setlength{\parskip}{6pt plus 2pt minus 1pt}}
}{% if KOMA class
  \KOMAoptions{parskip=half}}
\makeatother
\usepackage{graphicx}
\makeatletter
\newsavebox\pandoc@box
\newcommand*\pandocbounded[1]{% scales image to fit in text height/width
  \sbox\pandoc@box{#1}%
  \Gscale@div\@tempa{\textheight}{\dimexpr\ht\pandoc@box+\dp\pandoc@box\relax}%
  \Gscale@div\@tempb{\linewidth}{\wd\pandoc@box}%
  \ifdim\@tempb\p@<\@tempa\p@\let\@tempa\@tempb\fi% select the smaller of both
  \ifdim\@tempa\p@<\p@\scalebox{\@tempa}{\usebox\pandoc@box}%
  \else\usebox{\pandoc@box}%
  \fi%
}
% Set default figure placement to htbp
\def\fps@figure{htbp}
\makeatother
\setlength{\emergencystretch}{3em} % prevent overfull lines
\providecommand{\tightlist}{%
  \setlength{\itemsep}{0pt}\setlength{\parskip}{0pt}}
\usepackage{bookmark}
\IfFileExists{xurl.sty}{\usepackage{xurl}}{} % add URL line breaks if available
\urlstyle{same}
\hypersetup{
  pdftitle={ED\_Project\_Proposal Standard Deviants},
  pdfauthor={Ndumiso, Nyawo and 26396769; Lebohang, Maleke, 28293436; Masiu, Mohoang, 28298918; Rudairo, Nyamande, 28023366; Lwazi, Zwane, 28913884},
  hidelinks,
  pdfcreator={LaTeX via pandoc}}

\title{ED\_Project\_Proposal Standard Deviants}
\author{Ndumiso, Nyawo and 26396769 \and Lebohang, Maleke,
28293436 \and Masiu, Mohoang, 28298918 \and Rudairo, Nyamande,
28023366 \and Lwazi, Zwane, 28913884}
\date{Deadline: 31 July 15:00}

\begin{document}
\maketitle

\subsubsection{Instructions to using this
template}\label{instructions-to-using-this-template}

Delete this section when submitting the final document. You may not
alter the section headings below. That is, the document may only consist
of:

(1)Aim; (2) Experimental design. The final knitted PDF report may not
exceed 1 page (otherwise your group will get zero). Inside the header,
add the title ``ED\_Project\_Proposal '' (replace with your group name),
and place group member information at author. That is, names, surnames,
student numbers of members. The group leader should be places first, and
then the rest of the members should sorted alphabetically by surname.
E.g., Name1, Surname1, Student number1, Name2, Surname2, Student
number2, \ldots{} Read the instructions at the different sections
carefully. Delete the instruction text when submitting the final
document.

\#\#Briefly discuss the aim of the project. What is the research
question? Why is it important to answer this question?

\textbf{AIM} The aim of this study is to investigate whether a student's
faculty of study has an effect on the time required to solve a
standardised logical reasoning test. The study will compare the average
problem-solving times of students from different faculties using a
one-way Analysis of Variance (ANOVA).

\textbf{RESEARCH QUESTION} - Does faculty of study influence the average
time taken by Stellenbosch University students to solve a standardised
logical reasoning test?

\subsubsection{Experimental design}\label{experimental-design}

This is a very important section. Describe the experimental units,
treatments and measured responses. How many replicates are you planning
on including? Note that a minimum of 5 replicates are required for the
experiment.

\textbf{EXPERIMENTAL UNITS}

The experimental units are individual undergraduate students currently
enrolled at Stellenbosch University

\textbf{FACTOR (Independent Variable)}

Faculty of Study

\textbf{TREATMENT LEVELS} Engineering Science Economic and Management
Sciences (EMS)

\textbf{RESPONSE VARIABLES}

The total time (in seconds) required to complete a standardised logical
reasoning test.

\textbf{REPLICATION}

Ten students will be recruited from each faculty, giving a total sample
size of thirty(30) participants.

\end{document}
